% Paste this environment into the IEEEtran source.
% The preamble must already \input the sibling iscafig.tex.
\begin{figure}[t]
\centering
\begin{tikzpicture}[iscafig]
  \node[lbl, anchor=west] at (0,0.38) {\textbf{Tuple}};
  \rec{0}{+0}{fldimm}{3.55cm}{length}
  \rec{-0.52}{+4}{fldcore}{3.55cm}{element 0}
  \node[lbl, anchor=north west] at (0.1,-0.82) {\scriptsize elements follow inline};

  \begin{scope}[xshift=4.35cm]
    \node[lbl, anchor=west] at (0,0.38) {\textbf{Range}};
    \rec{0}{+0}{fldimm}{3.55cm}{start}
    \rec{-0.52}{+4}{fldimm}{3.55cm}{stop}
    \rec{-1.04}{+8}{fldimm}{3.55cm}{step}
  \end{scope}

  \begin{scope}[yshift=-2.55cm]
    \node[lbl, anchor=west] at (0,0.38) {\textbf{\texttt{ITER}} object, four words};
    \rec{0}{+0}{fldmem}{7.15cm}{sequence, or the current range value}
    \rec{-0.52}{+4}{fldcore}{7.15cm}{index, or range stop}
    \rec{-1.04}{+8}{fldcore}{7.15cm}{extra, or range step}
    \rec{-1.56}{+12}{fldimm}{7.15cm}{kind: 0 seq, 1 range, 2 dict, 3 set, 5 str}
    \node[lbl, anchor=west, align=left] at (0,-2.15)
      {List and range \texttt{FOR\_ITER} keeps this same four-word record in the ring,\\
       so the hot path does not reload the heap object on every step.};
  \end{scope}
\end{tikzpicture}
\caption{Iterables.
A tuple is a length word followed by its elements.
A range is three integers.
An \texttt{ITER} value points at a four-word record.
For a list the first word is the list and the second is the index.
For a range the first three words are the current value, the stop, and the step.
The kind in the last word selects the step.}
\label{fig:iter}
\end{figure}
